\section{Conclusiones}

La investigación teórica y el desarrollo práctico de la arquitectura BIGENIA
demuestran empíricamente la viabilidad de democratizar la analítica de datos
en el sector PyME mediante el uso de Inteligencia Artificial Generativa. La
implementación sobre un modelo de datos relacional de prueba confirma que la
orquestación de agentes logra abstraer la complejidad técnica subyacente,
consolidando una interfaz directa y efectiva entre el usuario de negocio y la
base de datos.

A través de la metodología de evaluación \textit{LLM-as-a-Judge}, aplicada
sobre un conjunto estructurado de 35 casos de prueba, el sistema evidenció un
cumplimiento robusto de los objetivos funcionales. El Agente de Intención
demostró una precisión perfecta (100\%) en la clasificación de solicitudes,
lo que subraya la resiliencia del sistema y la efectividad de sus mecanismos
de control (\textit{guardrails}) para identificar y mitigar consistentemente
las peticiones fuera de dominio. Por su parte, el Agente Visualizador logró
un desempeño sobresaliente con una tasa de éxito del 96.9\% en la generación
de representaciones gráficas adecuadas. El Agente de Base de Datos también
alcanzó niveles muy altos de efectividad (92.9\%) en la generación autónoma
de sentencias SQL precisas.

Si bien se identificó una degradación en la precisión de inferencia frente a
consultas que exigen múltiples uniones y subconsultas anidadas, este
comportamiento permite establecer con claridad el límite operativo de la
versión actual sin invalidar su utilidad estratégica. En definitiva, este
trabajo corrobora que el reemplazo de interfaces técnicas tradicionales por
modelos de procesamiento de lenguaje natural reduce drásticamente la curva de
aprendizaje y dota de autonomía analítica real a organizaciones con recursos de
TI limitados, sentando bases sólidas para el desarrollo de herramientas analíticas
de próxima generación.